%======================= CONCEPTION BASE DE DONNÉES =========================
\chapter{Conception de la base de données}

\section{Démarche}

Bien que le prototype fonctionne aujourd'hui sur des données \emph{mockées},
un \textbf{schéma de base de données cible} a été conçu et versionné, afin de
préparer le passage à une base \textbf{PostgreSQL} réelle sans réécrire les
routes. Ce schéma est décrit en détail dans le dépôt
(\texttt{docs/schema-base-de-donnees.md}) et matérialisé par des fichiers de
migration SQL (\texttt{database/migrations/}). Les modèles \textbf{SQLAlchemy}
du backend en constituent une première traduction.

\section{Entités logiques principales}

Le tableau~\ref{tab:entites} présente les entités logiques du domaine, leur
rôle et leur état de mise en \oe uvre actuel, en toute transparence.

\begin{table}[H]
\centering
\caption{Entités logiques et état actuel}
\label{tab:entites}
\small
% Première colonne élargie : les noms d'entités (ex. campaign_drafts) sont
% en police à chasse fixe et ne doivent pas déborder sur la colonne voisine.
\begin{tabularx}{\textwidth}{@{}p{3.1cm}L p{3.4cm}@{}}
\toprule
\textbf{Entité} & \textbf{Rôle} & \textbf{État actuel} \\
\midrule
\texttt{users} & Comptes de la plateforme (rôle et statut). & Mocké (en mémoire). \\
\texttt{advertisers} & Annonceurs clients, propriétaires des campagnes et magasins. & Mocké (en mémoire). \\
\texttt{campaigns} & Campagnes publicitaires drive-to-store. & Mocké (en mémoire). \\
\texttt{campaign\_drafts} & Brouillons produits par l'assistant de création. & \textbf{Réellement persisté} (PostgreSQL si disponible, sinon SQLite local). \\
\texttt{stores} & Magasins (points de vente) d'un annonceur. & Mocké (en mémoire). \\
\texttt{creatives} & Visuels / créatives d'une campagne. & Métadonnées en mémoire (aucun fichier stocké). \\
\texttt{statistics} & Statistiques de performance par campagne / magasin. & Mocké (agrégé uniquement pour le tableau de bord). \\
Paramètres de compte & Préférences du compte (devise, fuseau, langue, notifications). & En mémoire (deux registres distincts : par annonceur et global). \\
\bottomrule
\end{tabularx}
\end{table}

\section{Relations}

Les relations principales du schéma cible sont les suivantes :

\begin{itemize}
  \item un \textbf{annonceur} possède plusieurs \textbf{campagnes} et
        plusieurs \textbf{magasins} (relations 1--N) ;
  \item une \textbf{campagne} appartient à un seul annonceur, et porte
        plusieurs \textbf{créatives} et \textbf{statistiques} ;
  \item une \textbf{statistique} est rattachée à une campagne et,
        éventuellement, à un magasin (rattachement optionnel) ;
  \item un \textbf{brouillon de campagne} peut être rattaché à un annonceur
        (rattachement optionnel) ;
  \item la table \texttt{users} est, à ce stade, \textbf{indépendante} des
        autres (pas encore de clé étrangère vers un annonceur au niveau base).
\end{itemize}

\section{Choix de conception}

Le schéma cible retient des choix classiques et robustes : identifiants
auto-incrémentés, types numériques précis pour les montants, dates avec fuseau
horaire, contraintes de validation (valeurs autorisées, montants positifs,
cohérence des dates), et index sur les colonnes les plus filtrées (clés
étrangères, dates). Les brouillons de campagne stockent les champs principaux
en colonnes classiques et les sélections détaillées de l'assistant (appareils,
systèmes, formats, catégories) dans une colonne \textbf{JSON}, ce qui offre de
la souplesse sans migration à chaque évolution du formulaire.

\section{État honnête}

Il importe de rappeler que \textbf{PostgreSQL n'est pas branché en production}
aujourd'hui. Les modèles et migrations existent et sont cohérents, mais l'API
sert des données en mémoire, à l'exception notable des \textbf{brouillons de
campagne}, seuls réellement persistés. L'alignement complet des modèles sur le
schéma cible (par exemple l'ajout d'un mot de passe haché pour les
utilisateurs) fait partie des perspectives.
